\section{Introduction}
\label{sec:intro}

A \emph{symmetry} of a physical system is a transformation that maps solutions (or states) to
solutions and leaves the laws unchanged. Two kinds are familiar. \emph{Space-time} symmetries act on
the coordinates: rotations, boosts (together the Lorentz transformations) and translations.
\emph{Internal} symmetries act on the components of the fields at each point and commute with every
space-time transformation, like the isospin rotations of a doublet of fields or the phase rotations
whose generator is the electric charge. In the 1960s it was hoped that the two kinds could be united
in one larger symmetry, in which an internal multiplet would also be a Lorentz multiplet. A series of
theorems---McGlinn~\cite{McGlinn1964}, O'Raifeartaigh~\cite{ORaifeartaigh1965} and Coleman and
Mandula~\cite{ColemanMandula1967}---showed that this is impossible under hypotheses that seemed
harmless, and Haag, {\L}opusza\'nski and Sohnius~\cite{HLS1975} showed that the only way out within
operator algebras is to admit generators that \emph{anticommute}: supersymmetry. Pelc and Horwitz
extended the theorem of Coleman and Mandula to any number of space dimensions~\cite{PelcHorwitz1997}.

A theorem cannot have a counterexample in which all its hypotheses hold. What can exist is a
structure in which a hypothesis fails and the conclusion of the theorem fails with it: a
\emph{counterexample to the conclusion}. This paper constructs such structures by giving up one
hypothesis that every proof uses and none states as a physical assumption: \emph{associativity}.
The symmetries are assumed to compose like the elements of a group (associatively: three
transformations performed in succession give a result that does not depend on which two are combined
first), to act by linear operators, whose products are associative, and therefore to have generators
that obey the Jacobi identity. We show in Sec.~\ref{sec:nogo} exactly where these assumptions enter.

The non-associative object we use is the algebra of \emph{bi-octonions} $\mathbb B_{\rm LL}$, the
octonions with complex coefficients. Its unit elements form a \emph{Moufang loop}, a set with a
product that has a unit element and inverses but is not associative, which we call the
\emph{Lorentz super-loop} $\mathbb L_{\rm LL}$. It contains the group $SL(2,\mathbb C)$, the double cover of the
Lorentz group, and eight further directions, the \emph{spinor charges}, which transform as a Weyl
spinor of one chirality and form an isospin doublet with electric charges $\pm\tfrac12$. In its
algebra of infinitesimal elements, a Malcev algebra $\mathfrak m_{\rm LL}$, two spinor charges have a
Lorentz generator as their bracket. We prove (Theorem~L) that no Lie algebra can do this, that the
super-loop is the only solution among all Lorentz- and isospin-covariant brackets (Theorem~M), and
that its non-associativity is irreducible: it has no realization by operators (Theorem~R). These are
the requested counterexamples: Lorentz and internal degrees of freedom combined in one structure, with
operators whose eigenvalues are the electric charge, the chirality, the spinor number, and a new
quantity, the \emph{twist charge}, which takes the value $-2$ on the spinor charges and measures by how
much the Lorentz action on them fails to be a representation.

We then ask which manifolds carry the super-loop (Sec.~\ref{sec:manifolds}): the Minkowski spaces of
$7+1$ and $3+1$ dimensions inside the bi-octonions, the split octonions of signature $4+4$, the
seven-sphere and the super-loop itself, a parallelizable manifold that is not a Lie group. We analyze
the spin-zero superfluid of the spinor charges (Sec.~\ref{sec:superfluid}): the field content follows
from the bi-octonion product, the spin-zero pairs of one field are isospin triplets, the unitary state
wins, and the pair that is a singlet of the whole super-loop needs two species of fields. Finally we
decide whether these theories are renormalizable (Sec.~\ref{sec:renorm}).

Every statement of this paper is verified by a machine check, labelled as in Appendix~\ref{app:checks}
(for example A-2 or S-5); exact statements are checked in exact integer or rational arithmetic. The
programs, the data files and the commands that reproduce every number are listed in
Appendix~\ref{app:data}, and every symbol is defined in Appendix~\ref{app:notation}. Units: $\hbar$ and
the speed of light are 1; in Secs.~\ref{sec:superfluid} and~\ref{sec:renorm} the momentum cutoff
$\Lambda_{\rm LL}=1$ unless a unit is shown.

\begin{table}[t]
\caption{The theorems proved in this paper and the machine checks on which they rest.}
\label{tab:theorems}
\begin{ruledtabular}
\begin{tabular}{p{0.11\columnwidth}p{0.62\columnwidth}p{0.17\columnwidth}}
label & statement (short form) & checks\\
\hline
R & $\mathfrak m_{\rm LL}$ has no realization by operators; $\mathbb L_{\rm LL}$ is not a group of operators & A-6, A-7, A-8\\
L & no Lie algebra extends $\mathfrak{so}(3,1)$ by spinor charges whose brackets contain Lorentz generators & L-6, G-1, G-5, G-10\\
M & among all covariant brackets, the bi-octonion algebra $\mathfrak m_{\rm LL}$ is the unique such Malcev algebra; with translations or with isospin acting on the spinor charges there is none & G-3, G-4, G-7 to G-12, M-1 to M-9\\
SF1 & the spin-zero pairs of one field are isospin triplets & S-1, S-2\\
SF2 & the spectra of the superfluid & S-3\\
SF3 & the unitary state has the lowest energy & S-6\\
SF4 & the super-loop singlet pair needs two species & S-8, S-9\\
\end{tabular}
\end{ruledtabular}
\end{table}
